% 本文件由 LaTeX 排版 Agent 根据有效 translation.json 自动生成。
% author=Gregory Thaumaturgus; work_id=0608; title=论灵魂

\chapter{论灵魂}

% structure: p0001 source=h1 target=omitted reason=page-title-already-rendered-by-parent

% structure: p0002 source=h2 target=section reason=relative-heading-rank; base=h2

\section{序言}

最卓越的\Person{塔提安}，你曾嘱咐我们，为你撰写一篇关于灵魂的论述，并以有效的论证加以铺陈。你要求我们不必使用\WorkTitle{圣经}的见证——这方法本是向我们敞开的，且对于寻求虔敬心灵的人而言，比任何人的推理更能令人信服地阐述道理。然而，你说你渴望如此行，并非为了你自己的确信——因你已受教持守\WorkTitle{圣经}和传统，并避免因人的任何辩论巧思而动摇——而是为了驳斥那些持有不同见解、不承认应如此信赖\WorkTitle{圣经}，并试图凭借某种言语机巧来争取那些不谙此类讨论之人。因此，我们欣然接受这项委托，并未因不熟悉这种辩论方法而退缩，反倒因得知你对我们的善意而备受鼓舞。因为你对我们的友善与亲切，将使你懂得如何公开提出你所认可的我方正确表述，而对我方你认为不合宜的言论则予以略过和隐藏。深知此点，我便满怀信心着手阐释。在论述中，我将采用某种次序与连贯性，如同那些极为精通此事的人在向愿意明智探究任何主题的人所使用的那样。\par

首先，我将提议探究：灵魂按其本性可被认识的准则是什么；然后，灵魂存在的证明手段是什么；此后，它是\Emph{实体}还是\Emph{偶性}；接着，基于这些点，它是物体还是非物体；然后，它是单纯的还是复合的；接下来，它是可朽的还是不朽的；最后，它是有理性的还是无理性的。\par

这些问题正是通常在任何关于灵魂的探究中最为重要、最能界定其独特本性的讨论焦点。作为确立这些探究事项的论证，我们将采用那些通常能自然证明手头问题可信度的通用思考方式。但为求简洁与实用，我们目前仅使用对我们主题最具决定性的论证方式，以便清晰明了的观念能赋予我们应对反对者的一定准备。由此，我们开始讨论。\par

% structure: p0006 source=h2 target=section reason=relative-heading-rank; base=h2

\section{一、认识灵魂的准则何在}

一切存在的事物，要么被感官所认识，要么被思想所领会。而属于感官的事物，其充分的证明在于感官本身；因为一旦接触，它便在我们心中产生对底层事物的印象。但被思想所领会的事物，不是通过自身，而是通过其运作而被认识。因此，灵魂既然不通过自身而被认识，就应通过其效果而被恰当地认识。\par

% structure: p0008 source=h2 target=section reason=relative-heading-rank; base=h2

\section{二、灵魂是否存在}

我们的身体，当被推动时，或由外或由内被推动。它并非由外被推动，这从下面的事实可见：它既非通过推动，也非通过牵引而被推动，如同无灵魂之物。再者，如果它由内被推动，它并非像火那样按本性被推动。因为火只要存在，就不会失去其活动；而身体一旦死亡，便成为缺乏活动的身体。因此，如果它既不从外被推动（如同无灵魂之物），也不按本性被推动（如同火），那么显然它是被灵魂所推动，灵魂也赋予身体生命。那么，如果灵魂被证明赋予我们身体生命，灵魂也将通过其运作而被它自身所认识。\par

% structure: p0010 source=h2 target=section reason=relative-heading-rank; base=h2

\section{三、灵魂是否实体}

灵魂是实体，以下列方式得到证明。首先，因为赋予\Quote{实体}一词的定义非常适用于它。该定义的要旨是：实体是那种始终同一、就数目而言始终如一者，却能依次接纳相反者。而此灵魂，在不逾越其自身本性的界限的情况下，依次接纳相反者，我想这对每个人都是清楚的。因为正义与不义、勇敢与怯懦、节制与放纵，在灵魂中相继出现；这些是相反者。因此，如果实体的特性是能够依次接纳相反者，且灵魂被证明符合这一定义，那么灵魂便是实体。其次，因为如果身体是实体，灵魂也必须是实体。因为不能是：只被赋予生命的东西是实体，而赋予生命的东西不是实体——除非有人断言非存在是存在的原因；或者，除非有人疯狂到声称依赖物本身就是其所依赖之物（没有它便不能存在）的原因。\par

% structure: p0012 source=h2 target=section reason=relative-heading-rank; base=h2

\section{四、灵魂是否非物体}

灵魂在我们的身体中，这已在上文表明。因此，现在我们应当确定灵魂以何种方式存在于身体中。如果它与身体并置（如同小石子与另一小石子），那么灵魂将是物体，而且整个身体将不会全部被灵魂赋予生命，因为灵魂仅以某一部分与身体并置。但如果灵魂与身体混合或融合，那么灵魂将成为复合的，而非单纯的，从而被剥夺了灵魂所应具备的理性。因为复合的也是可分的、可分解的；而可分解的，反过来又是组合的；组合的则是可以三种方式分离的。此外，物体附于物体产生重量；但灵魂存在于身体中，不产生重量，反而赋予生命。因此，灵魂不能是物体，而是非物体的。\par

再者，如果\OriginalTerm{灵魂}{soul}是形体，那么它要么由外部推动，要么由内部推动。但它并非由外部推动；因为它不像无灵魂之物那样被推动或牵引。它也不是由内部推动，如同具有灵魂的对象；因为谈论灵魂的灵魂是荒谬的。因此，它不可能是一个形体，而是无形的。\par

此外，如果\OriginalTerm{灵魂}{soul}是形体，它就有可感觉的质，并靠滋养维持。但它并非这样被滋养。因为如果它被滋养，它不是像身体那样物质性地被滋养，而是无形地；因为它靠理性滋养。因此，它没有可感觉的质：因为正义、勇气或这些事物中的任何一个，都不是看得见的东西；然而这些是灵魂的质。因此，它不可能是一个形体，而是无形的。\par

更进一步，既然所有物质实体分为有生命的和无生命的，让那些主张\OriginalTerm{灵魂}{soul}是形体的人告诉我们，我们应当称它为有生命的还是无生命的。\par

最后，如果每一个形体都有颜色、大小和形状，而\OriginalTerm{灵魂}{soul}中无法察觉这些质中的任何一个，那就得出灵魂不是形体的结论。\par

% structure: p0018 source=h2 target=section reason=relative-heading-rank; base=h2

\section{五、灵魂是单一的还是复合的}

那么，我们证明\OriginalTerm{灵魂}{soul}是单一的，最好的方式是通过那些已经证明其无形性的论证。因为如果它不是形体，而每个形体都是复合的，复合的东西由部分构成，因而是多重的；相反，灵魂由于是无形的，就是单一的；因为它既非复合的，也非可分化为部分。\par

% structure: p0020 source=h2 target=section reason=relative-heading-rank; base=h2

\section{六、我们的灵魂是否不朽}

在我看来，必然的结论是，单一的便是\OriginalTerm{不朽}{immortal}。至于这如何得出，请听我的解释：没有任何存在物是它自身的败坏者，否则它从起初就不可能具有任何完全的一致性。因为易于朽坏的事物被对立面所败坏；因此，一切朽坏的事物都会分解；分解的东西是复合的；复合的东西由许多部分构成；由部分构成的东西显然由不同的部分构成；不同便不是同一；因此，灵魂既然是单一的，不由不同部分构成，而是非复合且不可分解的，那么基于此，它必然是不朽坏且不朽的。\par

再者，凡由他物推动、自身不拥有生命本原、而是从推动者得到生命本原的事物，只要被那在其中运作的能力所持有，就持续存在；当运作能力停止时，从它获得活动能力的事物也停止。但灵魂是\OriginalTerm{自动的}{self-acting}，它的存在不会停止。因为自动的事物是永恒活动的；永恒活动的是不间断的；不间断的是无终的；无终的是不朽坏的；不朽坏的是不朽的。因此，如果灵魂是自动的，如上所述，那么根据已陈明的推理方式，它是不朽坏且不朽的。\par

并且进一步，凡不被其自身固有的\OriginalTerm{恶}{evil}所败坏的事物，是不朽坏的；恶与善对立，因而是善的败坏者。因为身体的恶无非是痛苦、疾病和死亡；正如另一方面，它的卓越之处是美丽、生命、健康和活力。因此，如果灵魂不被其自身固有的恶所败坏，而灵魂的恶是怯懦、不节制、嫉妒等，而所有这些事物并不剥夺它生命和行动的能力，那么它便是不朽的。\par

% structure: p0024 source=h2 target=section reason=relative-heading-rank; base=h2

\section{七、我们的灵魂是否是理性的}

我们的灵魂是\OriginalTerm{理性的}{rational}，这一点可以用许多论证来证明。首先，从它发现了服务我们生活的技艺这一事实来看。因为没有人会说这些技艺是随意或偶然引入的，正如没有人能证明它们是空洞无用的，对我们的生活方式没有益处。因此，如果这些技艺有助于我们生活的有益之物，而有益之物是值得称赞的，值得称赞之物由理性构成，而这些都是灵魂的发现，那么我们的灵魂就是理性的。\par

再次，我们的灵魂是\OriginalTerm{理性的}{rational}，也通过我们感觉不足以认识事物这一事实得到证明。因为我们仅仅运用感觉机能并不足以认识事物。但我们不愿不经探究就停留在感觉中，这证明感觉离开了理性，被认为无法区分那些形式相同、颜色相似但本性截然不同的事物。因此，如果感觉离开了理性给我们关于事物的错误观念，我们就必须考虑事物本身是否能够被真正理解。如果它们能被理解，那么使我们能够达到它们的能力就不同于且优越于感觉。如果它们不被理解，那么我们根本不可能理解那些外表与现实不同的事物。但对象能被我们理解，这从以下事实可以清楚看出：我们以适于用途的方式使用每一事物，然后又随意转换它们。因此，如果已经证明事物能被我们理解，而感觉离开了理性是对对象的错误检验，那么由此得出，\OriginalTerm{理智}{intellect}是在理性中区分一切事物，并认识事物实际所是的。但理智正是灵魂的理性部分，因此灵魂是理性的。\par

最后，因为我们不做任何事先没有为自己确定的事情；而这无非就是灵魂的高超特权——因为对事物的知识并非从外部来到它那里，而是它仿佛用自己思想的装饰来摆出这些事物，从而先在自身中描绘出对象，然后才付诸实际行为——并且因为灵魂的高超特权无非就是凭\OriginalTerm{理性}{reason}做一切事，在这点上它也不同于感觉，灵魂由此已被证明是理性的。\par

\SourceRecord{Translated by S.D.F. Salmond.；Ante-Nicene Fathers；Vol. 6.；Edited by Alexander Roberts, James Donaldson, and A. Cleveland Coxe.；Buffalo, NY: Christian Literature Publishing Co.,；1886.；<http://www.newadvent.org/fathers/0608.htm>.}
